\documentclass[11pt,a4paper]{article}
\usepackage[utf8]{inputenc}
\usepackage[spanish]{babel}
\usepackage{graphicx,booktabs,hyperref,geometry}
\geometry{margin=2.5cm}
\title{Proyecto integrador: \emph{[título del proyecto]}\\
\large Redes de Sensores Inalámbricos -- Práctica 9}
\author{[Integrantes]}
\date{\today}
\begin{document}
\maketitle
\section{Escenario y requisitos}
Describa el caso de uso, el área a cubrir, el número de nodos, la frecuencia de muestreo requerida y la vida útil objetivo.
\section{Arquitectura}
Incluya un diagrama de bloques (nodos, border router, tunslip6, broker, panel) y la pila de protocolos de cada componente.
\section{Decisiones de diseño}
Capa MAC (CSMA/TSCH), periodo de envío, protocolo de aplicación (CoAP/MQTT), parámetros RPL. Justifique cada decisión con datos de las prácticas anteriores.
\section{Metodología experimental}
Simulaciones ejecutadas (tamaños, semillas, duración), métricas y scripts usados.
\section{Resultados}
Tablas y gráficas: PDR, saltos medios, latencia, consumo estimado y vida útil.
\section{Conclusiones y trabajo futuro}
\end{document}
